% concatenated from Frobenius_trace_DM_rk2.tex
\documentclass[20pt, oneside]{article}   	% use "amsart" instead of "article" for AMSLaTeX format
\usepackage{geometry}                		% See geometry.pdf to learn the layout options. There are 
\usepackage{amsmath}
\usepackage{amsthm}	
\usepackage{amssymb}
\usepackage{tikz}
\usepackage[implicit]{hyperref}
\usepackage[bottom]{footmisc}
\usepackage[affil-it]{authblk}
\usepackage[toc,page]{appendix}
\newtheorem{thm}{Theorem}[section]

\newtheorem{prop}[thm]{Proposition}
\newtheorem{note}[thm]{Note}
\newtheorem{lem}[thm]{Lemma}
\newtheorem{cor}[thm]{Corollary}
\newtheorem*{claim*}{Claim}
\newtheorem{lemma}{Lemma}
\newtheorem*{thm*}{Theorem}
\newtheorem*{corr}{Corollary}
\newtheorem*{theorem}{Theorem}
\newtheorem{thmm}{Main Theorem}
\newtheorem*{ques}{Question}
\newtheorem*{fact}{Fact}
\newtheorem{term}[thm]{Terminology}

\newtheorem{ex}[thm]{Example}
\newtheorem{man}{Theorem}

\theoremstyle{definition}
\newtheorem*{rem*}{Remark}
\newtheorem{defi}[thm]{Definition}
\newtheorem{rem}[thm]{Remark}
\newtheorem{conv}[thm]{Convention}



\newcommand{\F}{\mathbb{F}}
\newcommand{\fl}{\mathfrak{l}}
\newcommand{\fL}{\mathfrak{L}}
\newcommand{\fq}{\mathfrak{q}}
\newcommand{\fa}{\mathfrak{a}}
\newcommand{\fp}{\mathfrak{p}}
\newcommand{\Gal}{\mathrm{Gal}}
\newcommand{\Aut}{\mathrm{Aut}}
\newcommand{\sep}{\mathrm{sep}}
\newcommand{\alg}{\mathrm{alg}}
\newcommand{\GL}{\mathrm{GL}}
\newcommand{\cO}{\mathcal{O}}
\newcommand{\M}{\mathcal{M}}
\newcommand{\rsep}{r_{\rm sep}}

\newcommand{\hookuparrow}{\mathrel{\rotatebox[origin=c]{90}{$\hookrightarrow$}}}
\newcommand{\hookdownarrow}{\mathrel{\rotatebox[origin=c]{-90}{$\hookrightarrow$}}}
\newcommand{\xdownarrow}[1]{%
  {\left\downarrow\vbox to #1{}\right.\kern-\nulldelimiterspace}
}










\title{Frobenius Traces for Rank-2 Drinfeld Modules, Higher-Dimensional Galois Representations, and a Strong Multiplicity One Theorem in Positive Characteristic}
					% Activate to display a given date or no date
\author{Chien-Hua Chen \thanks{Electronic address: \texttt{cc45@aub.edu.lb}; ORCID: \texttt{0000-0003-3267-5603} ; Corresponding author}}
\affil{}

\begin{document}

%\date{}		
\maketitle
\abstract

In this paper, we prove that if the Frobenius traces agree at all but finitely many places, then two $l$-adic Galois representations, associated to rank-$2$ non-CM Drinfeld modules of generic characteristic, are isomorphic. As a generalization, we show that  the "Frobenius trace equality at all but finitely many places forces isomorphism" between two Galois representations over a local field of positive characteristic only holds under an absolute irreducibility assumption. Moreover, we formulate and prove a function field analogue of strong multiplicity one property for semisimple Galois representations over a local field of positive characteristic.

\section {Introduction}


In 1950's, the well-known Brauer-Nesbitt theorem implies that two semisimple Galois representations with image in $\GL_n$ over a field of characteristic $0$ are isomorphic if and only if the two representations have the same Frobenius trace for all but finitely many places. As an application to representations associated to elliptic curves, we can conclude that Frobenius traces can distinguish elliptic curves up to isogeny. In positive characteristic case, however, this principle becomes more subtle. The classical argument relies on Newton identities to reconstruct characteristic polynomials from traces, which requires divisibility conditions that might fail in characteristic $p$. As a result, Frobenius traces alone do not, in general, determine representations up to isomorphism. Understanding when such a determination remains possible then becomes a natural problem.

Our first main result focuses on the representations arising from Drinfeld modules. Under the specialization to $\fl$-adic Galois representations associated to Drinfeld modules, the Weil pairing gives a  relation between Frobenius determinant and Frobenius image of a representation associated to certain twisted Carlitz module. We then apply this relation to bypass the Newton identities to achieve the same goal as the elliptic curve case. 

\begin{thmm}[Theorem \ref{mainthm1}]
Let $q=p^e$ be a prime power, $A:=\F_q[T]$ with field of fraction $F:=\F_q(T)$. Let $\phi_1$, $\phi_2$ be two non-CM rank-$2$ Drinfeld modules over $K$, a finite extension of $F$, of $A$-characteristic $0$. Let $\fl$ be a prime ideal of $A$ and $\rho_{i,\fl}:\Gal(\bar{K}/K)\rightarrow \GL_2(F_\fl)$  be the $\fl$-adic Galois representations associated to $\phi_i$ for $i=1,2$. 

Suppose that the Frobenius traces ${\rm{tr}}(\rho_{i,\fl}({\rm Frob}_{v}))$ are equal for all but finitely many places $v$ of $K$, then $\phi_1$ and $\phi_2$ are $K$-isogenous Drinfeld modules. Consequently, $\rho_{1,\fl}\cong\rho_{2,\fl}$.
\end{thmm}

As a generalization of our Main Theorem 1, we show in Proposition \ref{Main1} that for continuous representation over local field of positive characteristic, ``Frobenius trace equality for all but finitely many places$\Rightarrow$ two representations are isomorphic'' still  works when one representation is absolutely irreducible and the other one is irreducible. Therefore, apply Proposition \ref{Main1} to $\fl$-adic Galois representation associated to rank-$r$ non-CM Drinfeld modules, we get the corollary below.

\begin{corr}[Corollary \ref{corr}]
Let $q=p^e$ be a prime power, $A:=\F_q[T]$ with field of fraction $F:=\F_q(T)$. Let $r\in \mathbb{Z}_{>0}$, and $\phi_1$, $\phi_2$ be two non-CM rank-$r$ Drinfeld modules over $K$, a finite extension of $F$, of $A$-characteristic $0$. Let $\fl$ be a prime ideal of $A$. Let $\rho_{i,\fl}:\Gal(\bar{K}/K)\rightarrow \GL_r(F_\fl)$  be the $\fl$-adic Galois representations associated to $\phi_i$ for $i=1,2$. 

Suppose that the Frobenius traces ${\rm{tr}}(\rho_{i,\fl}({\rm Frob}_{v}))$ are equal for all but finitely many places $v$ of $K$, then
$\phi_1$ and $\phi_2$ are $K$-isogenous and so $\rho_{1,\fl}\cong\rho_{2,\fl}$.


\end{corr}

Beyond the Brauer-Nesbitt theorem, there is a finer result by Rajan (Theorem 1\&2 in \cite{Ra98}) on what will happen if there are two semisimple Galois representations with image in $\GL_n$ over a field of characteristic $0$ whose Frobenius traces are equal for a positive Dirichlet density of places. To formulate a function field analogue of the result by Rajan, merely looking at Frobenius traces is not enough, hence we turn back to compare Frobenius characteristic polynomials between two representations. Our approach combines an algebraic form of Chebotarev density with a componentwise thinness condition on the Zariski closure of the image. Under this framework, we obtain criteria ensuring that two representations become isomorphic after restriction to a finite extension. Under additional hypotheses, the two representations differ only by a finite character twist.


\begin{thmm}[Theorem \ref{thmb}]
Let $K$ be a global function field, and  $\rho_1,\rho_2:G_K\rightarrow \GL_n(E)$ be two semisimple Galois representations over a local field $E$ of positive characteristic. Denote that
\begin{itemize}
\item $\rho:=(\rho_1,\rho_2):G_K\rightarrow \GL_r(E)\times\GL_r(E).$
\item$ H:=\overline{\rho(G_K)}^{\rm Zar}\subset \GL_r\times\GL_r$ is the Zariski closure of $\rho(G_K)$ in the affine algebraic group $\GL_r\times\GL_r$ over $E$.
\item $\Phi:=H/H^\circ$, where $H^\circ$ is the identity component of $H$.
\item $S:=\{v\in\Sigma_K\mid \rho\ {\rm unramified\ at\ }v\}$
\item$X:=\left\{(g_1,g_2)\in \GL_r\times\GL_r\mid \chi_{\rm char}(g_1)=\chi_{\rm char}(g_2) \right\}$, where $\chi_{\rm char}(g)$ is the characteristic polynomial of $g$. 
\item $\Sigma_X:=\{v\in S\mid \rho({\rm Frob}_v)\in X\}$
\end{itemize}

Suppose $H^\circ\subset X$, then we have the followings:
\begin{enumerate}
\item[(i)] There is a finite extension $L/K$ such that $\rho_1|_{G_L}\cong\rho_2|_{G_L}$
\item[(ii)]  If $\rho_1$ is absolutely irreducible and $\overline{\rho_1(G_K)}^{\rm Zar}\subset \GL_r$ is connected, then there is a finite character $\chi:G_K\rightarrow E^*$ such that $\rho_2\cong\rho_1\otimes \chi$.
\end{enumerate}
\end{thmm}

Besides, we note in Remark \ref{rem1} the difference between the positive characteristic case and Rajan's result (Theorem 2) in \cite{Ra98}. In Remark \ref{rem2}, we apply our strong multiplicity one property for the situation where the two representations $\rho_1, \rho_2$ are $\fl$-adic Galois representations associated to non-CM Drinfeld modules of generic characteristic. 

\section*{Acknowledgement}

The main topic of this paper was originally raised by Professor Mihran Papikian to the author. He would like to thank Professor Mihran Papikian  for his helpful and inspiring discussions.





\section{Preliminaries}

\subsection{Brauer-Nesbitt Theorem and its related results}


\begin{thm}[Brauer-Nesnbitt]\label{BN}

Let $G$ be a group and $E$ be a field. If $\rho_1, \rho_2: G\rightarrow {\rm GL}_n(E)$ are two finite dimensional semisimple representations such that the characteristic polynomial of $\rho_1(g)$ and $\rho_2(g)$ coincide for all $g\in G$. Then $\rho_1$ and $\rho_2$ are isomorphic representations. 

\end{thm}



Note that the Brauer-Nesbitt theorem is applied to representations in arbitrary characteristic. However, the following well known application is only restricted to the characteristic $0$ case.

\begin{thm}[ \cite{Lan84}, Corollary 3.8]\label{0case}
Let $G$ be a group and $E$ be a field of characteristic $0$. If $\rho_1, \rho_2: G\rightarrow {\rm GL}_n(E)$ are two finite dimensional semisimple representations such that the traces ${\rm tr}(\rho_1(g))$ and ${\rm tr}(\rho_2(g))$ coincide for all $g\in G$. Then $\rho_1$ and $\rho_2$ are isomorphic representations. 

\end{thm}
For positive characteristic case, the statement above hold under an additional assumption that ${\rm char}(E)>n$:

\begin{prop}\label{pcase}
Let $G$ be a group $E$ be a field of characteristic $p$ with $p>n$. If $\rho_1, \rho_2: G\rightarrow {\rm GL}_n(E)$ are two finite dimensional semisimple representations such that the traces ${\rm tr}(\rho_1(g))$ and ${\rm tr}(\rho_2(g))$ coincide for all $g\in G$. Then $\rho_1$ and $\rho_2$ are isomorphic representations. \end{prop}

\begin{proof}[sketch of proof]
The key strategy to proof of Proposition \ref{pcase}  is to use Newton identities to deduce that one can rebuild characteristic polynomial of $\rho(g)$ from trace of  powers ${\rm tr}(\rho(g)^k)$. Suppose $\rho_1(g)$ has eigenvalues $\alpha_1,\cdots,\alpha_n$. Set power sums $p_k={\rm tr}(\rho_1(g)^k)$ for $k\geqslant 1$, and symmetric sums $e_k=\Sigma_{1\leqslant i_1\leqslant\cdots\leqslant i_k\leqslant n}\alpha_{i_1}\cdots \alpha_{i_k}$ with $e_0=1$. Let us write down the characteristic polynomial of $\rho_1(g)$ as $P_v(X)=X^n-e_1X^{n-1}+e_2X^{n-2}+\cdots+(-1)^ne_n$.
Since ${\rm char}(E)=p>n$, we have $k$ is invertible in $E$ for $1\leqslant k\leqslant n$. From the Newton identities $$ke_k=\Sigma_{i=1}^k(-1)^{i-1}e_{k-i}p_i \textrm{ for $1\leqslant k\leqslant n$ with $e_0=1$},$$ one can conclude that each $e_k$ is completely determined by $p_1, p_2,\cdots, p_k$, hence by ${\rm tr}(\rho_1(g^i)$ for $1\leqslant i\leqslant k$. Therefore, the characteristic polynomial of $\rho_1(g)$ and $\rho_2(g)$ are the same for all $g\in G$, the Brauer-Nesbitt theorem applies.
\end{proof}

Consider the case when $G=G_K$ is the absolute Galois group of a global field $K$. Since the trace function is a continuous class function, and Frobenius elements are dense in $G_K$ by Chebotarev density theorem, it is enough to test on trace of Frobenius ${\rm tr}(\rho_i({\rm Frob}_v))$ for all but finitely many places $v$ of $K$.



In 1998, Rajan studied strong multiplicity one for $\ell$-adic representations, and provided a finer result on representations having a positive density of places with equal Frobenius trace.  

\begin{thm}[\cite{Ra98}, Theorem 1] \label{Ra1}
Let $K$ be a global field , $G_K:=\Gal(\bar{K}/K)$ be the absolute Galois group, and $\Sigma_K$ be the set of non-archimedean places of $K$. Let $F$ be a non-archimedean local field of characteristic $0$ and residue characteristic $\fl$. Consider representations $\rho: G_K\rightarrow \GL_r(F)$ such that $\rho$ is continuous, unramified outside a finite set $S\subset\Sigma_K$, and semisimple.

 Suppose $\rho_1, \rho_2$ are such representations and the upper density of ${\rm SM}(\rho_1,\rho_2):=\{v\in\Sigma_K-S\mid {\rm tr}(\rho_1({\rm Frob}_v))={\rm tr}(\rho_2({\rm Frob}_v))\}$ in $\Sigma_K$ is strictly greater than $1-1/2r^2$, then $\rho_1\simeq\rho_2$.

\end{thm}

As an application to his main result, the following corollary on representations associated to elliptic curves gives a finer result on Theorem \ref{0case}. It shows what will occur when two $\ell$-adic representations associated to non-CM elliptic curves having equal Frobenius traces for positive density of places.

\begin{cor}[\cite{Ra98}, Corollary 2]\label{Ra2}
Suppose $E_1, E_2$ are two elliptic curves defined over a number field $K$. Let $S$ be the finite set consisting of the ramified places of $E_1$ and $E_2$. For a rational prime $\fl$, let $\rho_i$ be the $\fl$-adic Galois representation associated to $E_i$ for $i=1,2$. For $v\in\Sigma_K-S$ with $(\fl,v)=1$, let $a_v(E_i)$ denote the Frobenius trace ${\rm tr}(\rho_i({\rm Frob}_v))$ associated to $\rho_i$.  

Suppose that the set $\{v\in\Sigma_K-S\mid a_v(E_1)=a_v(E_2)\}$ in $\Sigma_K$ has positive upper density. Then there exists a quadratic Dirichlet character $\chi$ of $G_K$ such that $E_2$ is isogenous to $E_1\otimes \chi$.



\end{cor}




\subsection{Drinfeld modules and associated Galois representations}
Let $A=\F_q[T]$ be the polynomial ring over finite field with $q=p^e$ an odd prime power, $F=\F_q(T)$ be the fractional field of $A$, and  $K$ be a finite extension over $F$. Set $K\{\tau\}$ to be the twisted polynomial ring with usual addition rule, and the multiplication rule is defined to be $\tau\alpha=\alpha^q\tau \text{ for any } \alpha\in K$.  

We view $K$ as an $A$-field, which is a field equipped with a homomorphism $\gamma: A\rightarrow K$. The {\bf{$A$-characteristic}} of $K$ is defined to be the kernel of $\gamma$.


\begin{defi}
\begin{enumerate}
\item[(i)] Let $\phi$ be a Drinfeld $A$-module over $K$. The Drinfeld module $\phi$ gives $K^{{\rm{alg}}}$ an $A$-module structure, where $a\in A$ acts on $\alpha\in K^{{\rm{alg}}}$ via $a\cdot\alpha=\phi_{a}(\alpha)$. We use the notation $^{\phi}K^{{\rm{alg}}}$ to emphasize the action of $A$ on $K^{{\rm{alg}}}$.


\item[(ii)] Let $\mathfrak{a}$ be an ideal of $A$ with monic generator $a$.The {{$\mathfrak{a}$-torsion}} of $\phi$ is defined to be the set $$\phi[\mathfrak{a}]=\phi[a]:=\{ \alpha\in K^{\rm alg}\mid \phi_a(\alpha)=0 \}\subseteq K^{\rm alg}.$$ The action $b\cdot \alpha=\phi_b(\alpha) \ \forall\ b\in A, \forall \alpha\in\phi[\mathfrak{a}]$ also gives $\phi[\mathfrak{a}]$ an $A$-module structure. 
\end{enumerate}
\end{defi}

\begin{prop}\label{prop0.2}
Let $\phi$ be a rank $r$ Drinfeld module over $K$ and $\mathfrak{a}$ a non-zero ideal of $A$,
 If $\phi$ has $A$-characteristic prime to $\mathfrak{a}$, then the $A/\mathfrak{a}$-module $\phi[\mathfrak{a}]$ is free of rank $r$
\end{prop}
\begin{proof}
See \cite{G96} Proposition 4.5.7.
\end{proof}

Let $\phi$ be a rank $r$ Drinfeld module over $K$ of generic characteristic, then $\phi_a(x)$ is separable, so we have $\phi[a]\subseteq K^{{\rm{sep}}}$.
This implies that $\phi[a]$ has a $G_K$-module structure. Given a nonzero prime ideal $\fl$ of $A$ with prime power $\fl^i$, we can consider the $G_K$-module $\phi[\mathfrak{l}^i]$. We obtain the residual representation 
$$\bar{\rho}_{\phi,\mathfrak{l}^i}:G_K\longrightarrow {\rm{Aut}}(\phi[\mathfrak{l}^i])\cong GL_r(A/{\mathfrak{l}}^i).$$
In the case $i=1$, we get the so-called \textbf{mod $\mathfrak{l}$ Galois representation}
$$\bar{\rho}_{\phi,\mathfrak{l}}:G_K\longrightarrow {\rm{Aut}}(\phi[\mathfrak{l}])\cong GL_r(A/\fl).$$
Taking inverse limit with respect to $\mathfrak{l}^i$, we have the \textbf{$\mathfrak{l}$-adic Galois representation}



\section{Frobenius Traces and Isomorphism of Galois Representations}
As a alternative approach to Proposition \ref{pcase}, we observe that if one restrict the representations raised from rank-$2$ Drinfeld modules, one can still rely on Frobenius traces to distinguish Drinfeld modules up to isogeny classes. This method avoids the assumption ${\rm char}(E)>n$, while it requires a closed formula between coefficients of Drinfeld modules and coefficients of Frobenius characteristic polynomials, which in rank $r\geqslant 3$ is not clear.



\begin{lem}\label{existh}

There is a finite Galois extension $L/K$, independent of choice of $\fl$, such that $\rho_{1,\fl}\cong\rho_{2,\fl}$ when restricting to the open normal subgroup $H:=\Gal(\bar{K}/L)$ of $G_K$.

\end{lem}
\begin{proof}
Let us set $\phi_{1,T}=T+g_1\tau+g_2\tau^2$, and $\phi_{2,T}=T+s_1\tau+s_2\tau^2$, one can choose $L$ to be the compositum of Kummer extensions $L=K(\sqrt[q-1]{g_2},\sqrt[q-1]{s_2})$.
As we know that ${\rm tr}(\rho_{1, \fl}({\rm Frob}_v))={\rm tr}(\rho_{2, \fl}({\rm Frob}_v))\textrm{ for all but finitely many places v}$, Chebotarev density theorem then implies ${\rm tr}(\rho_{1,\fl}(g))={\rm tr}(\rho_{2,\fl}(g))$ for any $g\in G_K$, hence $${\rm tr}(\rho_{1,\fl}(h))={\rm tr}(\rho_{2,\fl}(h)) \textrm{ for any $h\in H$}.$$
 On the other hand, from the Weil paring for Drinfeld modules, we have that $$\det\circ\rho_{i,\fl}=\rho_{\psi_i,\fl} \textrm{ for $i=1,2$},$$ where $\psi_{1,T}=T-g_2\tau$, and $\psi_{2,T}=T-s_2\tau$ are twisted Carlitz module defined over $K$ together with their associated $\fl$-adic Galois representations $\rho_{\psi_i,\fl}:H\rightarrow \GL_2(F_\fl)$. Consider $\psi_1$ and $\psi_2$ over $L$, both of them are isomorphic to the twisted Carlitz module $\psi_T=T-\tau$. Therefore, one can deduce that $\rho_{\psi_1,\fl}\cong\rho_{\psi_2,\fl}$, which implies that $$\det\circ(\rho_{1,\fl}(h))=\det\circ(\rho_{2,\fl}(h)) \textrm{ for any $h\in H$}.$$ 
Thus Theorem \ref{BN} implies $\rho_1\cong\rho_2$ when restricting to the subgroup $H=\Gal(\bar{K}/L)$.

\end{proof}


\begin{lem}\label{1coho}

Let $G$ be a profinite group, and $H\leqslant G$ is a open normal subgroup with finite quotient $G/H$. Let $\rho_1, \rho_2: G\rightarrow \GL(V)\cong \GL_n(E)$ are surjective representations over a field $E$ that satisfy the properties below:
\begin{enumerate}
\item $\rho_i$ restricted to $H$ is irreducible for $i=1,2$.
\item There exists an $H$-equivariant isomorphism $\Phi:V\rightarrow V$ such that $\Phi\circ\rho_1(h)=\rho_2(h)\circ\Phi\textrm{ for all $h\in H$}$. 
\item $\rho_1(G)$ is open in $\mathrm{GL}_n(E)$.
\end{enumerate}
Then either $\rho_1\cong\rho_2$ or there exists a character $\chi:G\rightarrow E^*$ such that $\rho_1\cong\rho_2\otimes \chi$.





\end{lem}

\begin{proof}
Suppose that $\rho_1$ and $\rho_2$ are not isomorphic. Define $$\Psi:G\rightarrow {\rm Aut}_H(V)\textrm{ via } g\mapsto \Phi\rho_1(g)\Phi^{-1}\rho_2(g)^{-1}.$$

One then check that $\Psi(g_1g_2)=\Psi(g_1)\left( \rho_2(g_1)\Psi(g_2)\rho_2(g_1^{-1}) \right)$, so $\Psi(g)$ forms a $1$-cocycle, and this defines a nontrivial class in $H^1(G/H,{\rm Aut}_H(V))$. 

Now we consider the centralizer of the image  $\rho_1(H)$ in ${\rm End}_H(V)$, i.e. the set $$Z_{{\rm End}_E(V)}(\rho_1(H)):=\{f\in {\rm End}_E(V)\mid f\rho_1(h)=\rho_1(h)f \textrm{ for all }h\in H\}.$$


Since $\rho_1(G)$ is open in $\mathrm{GL}_n(E)$, it is Zariski dense in $\mathrm{GL}_n$. Since $H$ has finite index in $G$, the subgroup $\rho_1(H)$ has finite index in $\rho_1(G)$, and hence is also Zariski dense in $\mathrm{GL}_n$. Therefore
\[
Z_{\operatorname{End}_E(V)}(\rho_1(H))
=
Z_{\operatorname{End}_E(V)}(\mathrm{GL}_n)
=
E.
\]

 We can conclude that ${\rm Aut}_H(V)=E^*$, so any one class in $H^1(G/H,{\rm Aut}_H(V))$ corresponds to a character $\chi:G\rightarrow E^*$ that factors through the finite quotient $G/H$. From the equality $\Psi(g)=\Phi\rho_1(g)\Phi^{-1}\rho_2(g)^{-1}$ we get $$\Phi\rho_1(g)\Phi^{-1}=\Psi(g)\rho_2(g)=\chi(g)\rho_2(g) \textrm{ for some character $\chi$}.$$
 Thus we have $\rho_1\cong\rho_2\otimes\chi$.
 
\end{proof}





\begin{thm}\label{mainthm1}
Let $q=p^e$ be a prime power, $A:=\F_q[T]$ with field of fraction $F:=\F_q(T)$. Let $\phi_1$, $\phi_2$ be two non-CM rank-$2$ Drinfeld modules over $K$, a finite extension of $F$, of $A$-characteristic $0$. Let $\fl$ be a prime ideal of $A$ and $\rho_{i,\fl}:\Gal(\bar{K}/K)\rightarrow \GL_2(F_\fl)$  be the $\fl$-adic Galois representations associated to $\phi_i$ for $i=1,2$. 

Suppose that the Frobenius traces ${\rm{tr}}(\rho_{i,\fl}({\rm Frob}_{v}))$ are equal for all but finitely many places $v$ of $K$, then $\phi_1$ and $\phi_2$ are $K$-isogenous Drinfeld modules. Consequently, $\rho_{1,\fl}\cong\rho_{2,\fl}$.
\end{thm}


\begin{proof}
For the given two non-CM Drinfeld modules $\phi_1$ and $\phi_2$, from the Open Image Theorem and Lemma \ref{existh}, one can choose a prime $\tilde{\fl}\in A$ of large enough degree such that the following conditions are satisfied:
\begin{enumerate}
\item $\rho_{i,\tilde{\fl}}|_H:H\rightarrow \GL_2(F_{\tilde{\fl}})$ is irreducible for $i=1,2$
\item ${\rm Im}\rho_{i,\tilde{\fl}}\cong \GL_2(F_{\tilde{\fl}})$ for $i=1,2$
\item $\rho_{1,\tilde{\fl}}\cong\rho_{2,\tilde{\fl}}$ when restricted on $H=\Gal(\bar{K}/L)$
\end{enumerate}

Lemma \ref{1coho} then implies either $\rho_{1,\tilde{\fl}}\cong\rho_{2,\tilde{\fl}}$ or there is a nontrivial Galois character $\chi:G\rightarrow F_{\tilde{\fl}}^*$ such that $\rho_{1,\tilde{\fl}}\cong\rho_{2,\tilde{\fl}}\otimes\chi$. 

On the other hand, From the assumption that for all but finitely many places $v$ of $K$, we have ${\rm tr}(\rho_{1, \tilde{\fl}}({\rm Frob}_v))={\rm tr}(\rho_{2, \tilde{\fl}}({\rm Frob}_v))$. This forces $\rho_{1,\tilde{\fl}}\cong\rho_{2,\tilde{\fl}}$, otherwise there exists a nontrivial Galois character $\chi$ such that ${\rm tr}(\rho_{1, \tilde{\fl}}({\rm Frob}_v))={\rm tr}(\rho_{2, \tilde{\fl}}({\rm Frob}_v))\cdot \chi({\rm Frob}_v)$ for all but finitely many places $v$ of $K$, which forces ${\rm tr}(\rho_{1, \tilde{\fl}}({\rm Frob}_v))={\rm tr}(\rho_{2, \tilde{\fl}}({\rm Frob}_v))=0$ for all but finitely many places $v$, a contradiction.


Once we obtain an isomorphism $\rho_{1,\tilde{\fl}}\cong\rho_{2,\tilde{\fl}}$, we have immediately that there is a $G_K$-equivariant isomorphism between the Tate modules of $\phi_1$ and $\phi_2$, i.e. we have $V_{\tilde{\fl}}(\phi_1)\simeq V_{\tilde{\fl}}(\phi_2)$ where $V_{\tilde{\fl}}(\phi_i)=T_{\tilde{\fl}}(\phi_i)\otimes F_{\tilde{\fl}}$. Now Tate's isogeny theorem for Drinfeld modules gives
\[
\operatorname{Hom}_K(\phi_1,\phi_2)\otimes_A F_{\tilde{\fl}}
\simeq
\operatorname{Hom}_{G_K}(V_{\tilde{\fl}}(\phi_1),V_{\tilde{\fl}}(\phi_2)).
\]
The right-hand side is nonzero, since it contains an isomorphism. Hence $\operatorname{Hom}_K(\phi_1,\phi_2)\neq 0.$ Since \(\phi_1\) and \(\phi_2\) have the same rank, any nonzero homomorphism between them is an isogeny. Thus \(\phi_1\) and \(\phi_2\) are \(K\)-isogenous and so $\rho_1\cong\rho_2$.



\end{proof}




The key strategy in Theorem \ref{mainthm1} is that one can take a suitable finite extension $L/K$ to make Frobenius determinant to be equal, then the fact that Frobenius characteristic polynomial only involves in trace and determinant forces Frobenius characteristic polynomial of two representations are the same when restrict to $G_L$. Hence this technique would not work for higher dimensional Galois representations even for representations associated to rank-$r$ Drinfeld modules with $r\geqslant 3$. In order to generalize Theorem \ref{mainthm1} to arbitrary $n$-dimensional semisimple $\fl$-adic Galois representations over a local field $E$ of positive characteristic, we further requires a condition on irreducibility of two representations.


\begin{prop}\label{Main1}
Let $K$ be a global function field, $G_K:=\Gal(\bar{K}/K)$, and $E$ be any local field of characteristic $p>0$. Let $n\in \mathbb{N}$, and $\rho_1, \rho_2:G_K\rightarrow \GL_n(E)$ be continuous representations, unramified outside a finite set $S$ of places of $K$. Suppose further that
\begin{itemize}
\item For any $ \fp\not\in S$, we have ${\rm Tr}(\rho_1({\rm Frob}_\fp))={\rm Tr}(\rho_2({\rm Frob}_\fp))$.
\item $\rho_1$ is absolutely irreducible and $\rho_2$ is irreducible.
\end{itemize}
Then $\rho_1\cong\rho_2$
\end{prop}

\begin{proof}

We start with the fact that for $i=1,2$,  the composition $${\rm Tr}\circ \rho_i: G_K\rightarrow E\ \ \ g\mapsto {\rm Tr}(\rho_i(g))$$ is
continuous and invariant under conjugation. Hence Chebotarev density Theorem implies that ${\rm Tr}(\rho_1(g))={\rm Tr}(\rho_2(g))$ for any $g\in G_K$.




Now we consider the well-known Burnside's theorem below:
\begin{lem}[Burnside's Theorem]
Let $E$ be a field, $V$ be a finite dimensional vector space over $E$, $K$ be a  field with $G_K:=\Gal(\bar{K}/K)$, and $\rho:G_K\rightarrow GL_E(V)$ be a representation. Set $A:=E[\rho(G_K)]\subset {\rm End}_E(V)$. If $\rho$ is absolutely irreducible, then $A={\rm End}_E(V)$.
\end{lem}

Set $A_1=E[\rho_1(G_K)]={\rm End}_E(V_1)$ and $A_2=E[\rho_2(G_K)]$, where $V_1$ is a $n$-dimensional vector space over $E$ with $G_K$-action defined via $\rho_1$, similar for the vector space $V_2$. We fix a $E$-basis $\{a_i\}_{i=1}^{n^2}$ of $A$, where $a_i=\rho_1(g_i)$ for some $g_i\in G_K$. Let $\{a_i^*\}_{i=1}^{n^2}$ be the trace-dual basis of ${\rm End}_E(V_1)$, i.e. ${\rm Tr}(a_i^*a_j)=\delta_{ij}$. Then we define the $E$-linear map

$$
\begin{array}{lrl}
\Pi:&{\rm Hom}_E(V_1,V_2)&\rightarrow {\rm Hom}_E(V_1,V_2)\\
&S& \mapsto \sum_{i=1}^{n^2}\rho_2(g_i)Sa_i^*
\end{array}
$$
On the other hand, we define the following $E$-algebra homomorphism

$$\Phi:{\rm End}_E(V_1)\rightarrow E[\rho_2(G_K)]\ \ \sum_{i=1}^{n^2} c_ia_i\mapsto \sum_{i=1}^{n^2}c_i\rho_2(g_i).$$

Note that $\Phi$ is well-defined since for $\sum_{i=1}^{n^2}c_ia_i=0$, we have  ${\rm Tr}\left(\left [\sum_{i=1}^{n^2}c_i \rho_1(g_i)\right ]\rho_1(g)\right)=0$ for any $g\in G_K$. Its image via $\Phi$ has trace 
$${\rm Tr}\left(\left [\sum_{i=1}^{n^2}c_i \rho_2(g_i)\right ]\rho_2(g)\right)=0 \textrm{ for all $g\in G_K$}.$$ 
As $\rho_2$ is semisimple, $A_2$ is a semisimple algebra, thus the trace pairing on $A_2$ is non-degenerate, therefore the above trace equality implies $\sum_{i=1}^{n^2}c_i \rho_2(g_i)=0$. 



Now we claim that for $S\in {\rm Hom}_E(V_1,V_2)$, $\Pi(S):V_1\rightarrow V_2$ is $G_K$-equivariant, and $\Pi$ is not a zero map. Given a fixed $g\in G_K$, we can write 
\begin{equation}
\rho_1(g) a_i = \sum_{j=1}^{n^2} \text{Tr}(a_j^* \rho_1(g) a_i) a_j \tag{1}, \textrm{ and set }c_{ij}(g)=\text{Tr}(a_j^* \rho_1(g) a_i).\end{equation}

Similarly, we have \begin{equation}
a_j^* \rho_1(g) = \sum_{i=1}^{n^2} \text{Tr}(a_j^* \rho_1(g) a_i) a_i^* \tag{2}, \textrm{ and set }d_{ji}=\text{Tr}(a_j^* \rho_1(g) a_i).\end{equation}
One can easily see that $c_{ij}(g)=d_{ji}(g)$. Now we compute
$$
\begin{array}{ll}
\rho_2(g)\Pi(S)&=\sum_{i=1}^{n^2}\rho_2(gg_i)Sa_i^*=\sum_{i=1}^{n^2}\Phi(\rho_1(g)a_i)Sa_i^*=\sum_{i=1}^{n^2}\sum_{j=1}^{n^2}c_{ij}(g)\Phi(a_j)Sa_i^*\\
&\ \\
&=\sum_{i=1}^{n^2}\sum_{j=1}^{n^2}c_{ij}(g)\rho_2(g_j)Sa_i^*=\sum_{j=1}^{n^2}\rho_2(g_j)S\left(\sum_{i=1}^{n^2} d_{ji}(g)a_i^*\right)\\
&\ \\
&=\sum_{j=1}^{n^2}\rho_2(g_j)Sa_j^*\rho_1(g)=\Pi(S)\rho_1(g).
\end{array}
$$
It remains to prove $\Pi$ is not a zero map, we start with the operator $\Pi(S)=\sum_i \rho_2(g_i)\, S\, a_i^\ast.$
Consider the equality $\Phi(a_i)=\rho_2(g_i)$. We have seen that $\Phi:{\rm End}_E(V_1)\rightarrow E[\rho_2(G_K)]\subset {\rm End}_E(V_2)$ is an $E$-algebra homomorphism. As ${\rm End}_E(V_1)$ is a simple algebra, $\Phi$ must be injective. Comparing $E$-dimension on ${\rm End}_E(V_i)$ for $i=1,2$ forces $\Phi$ to be an $E$-algebra isomorphism. Thus apply Skolem--Noether theorem we may write
\[
\Phi(a_i)=Ta_iT^{-1}
\]
for some $E$-linear isomorphism $T:V_1 \to V_2$. Hence we obtain
\[
\Pi(S)=\sum_i T a_i T^{-1} S a_i^\ast.
\]
Now choose $S = TX$ for some $X \in {\rm End}_E(V_1)$. Then
\[
\Pi(TX)
=
\sum_i T a_i T^{-1} T X a_i^\ast
=
T\left(\sum_i a_i X a_i^\ast\right).
\]
Using the identity
$$\sum_i a_i X a_i^\ast = {\rm Tr}(X)\,\mathrm{Id},$$
we get $\Pi(TX)={\rm Tr}(X)\,T.$ Put $X=E_{11}$ to be the elementary matrix, we have ${\rm Tr}(E_{11})=1$, and hence
\[
\Pi(T E_{11}) = T \neq 0.
\]
This proves the claim.

Finally, we choose a  non-zero map $S$ such that $\Pi(S)\neq0$, then apply Schur's lemma on the nonzero $G_K$-module homomorphism $\Pi(S):V_1\rightarrow V_2$ to get an isomorphism between $V_1$ and $V_2$, hence $\rho_1\cong\rho_2$.



\end{proof}

\begin{rem}
The assumption that $\rho_1$ in Proposition \ref{Main1} to be absolutely irreducible is necessary because we need it to apply Burnside's Theorem; to the author's knowledge, there is no counter-example showing the existence of two irreducible representation $\rho_1,\rho_2:G_K\rightarrow \GL_n(E)$ such that they have the same Frobenius trace for all but finitely many places while $\rho_1\not\cong\rho_2$.


\end{rem}

The below corollary is an application of Proposition \ref{Main1} on $\fl$-adic Galois representations arising from generic characteristic Drinfeld modules of arbitrary rank:

\begin{cor}\label{corr}
Let $q=p^e$ be a prime power, $A:=\F_q[T]$ with field of fraction $F:=\F_q(T)$. Let $r\in \mathbb{Z}_{>0}$, and $\phi_1$, $\phi_2$ be two non-CM rank-$r$ Drinfeld modules over $K$, a finite extension of $F$, of $A$-characteristic $0$. Let $\fl$ be a prime ideal of $A$. Let $\rho_{i,\fl}:\Gal(\bar{K}/K)\rightarrow \GL_r(F_\fl)$  be the $\fl$-adic Galois representations associated to $\phi_i$ for $i=1,2$. 

Suppose that the Frobenius traces ${\rm{tr}}(\rho_{i,\fl}({\rm Frob}_{v}))$ are equal for all but finitely many places $v$ of $K$, then
$\phi_1$ and $\phi_2$ are $K$-isogenous and so $\rho_{1,\fl}\cong\rho_{2,\fl}$.

\end{cor}

\begin{proof}
 
From the fact that Frobenius traces are independent of the choice of prime $\fl$, and the  open image theorem for non-CM Drinfeld modules in \cite{PR09}, one can  deduce that there is a prime $\tilde{\fl}$ of sufficiently large degree so that that $\rho_{1,\tilde{\fl}}$ is absolutely irreducible and $\rho_{2,\tilde{\fl}}$ is irreducible. Thus we have $\rho_{1,\tilde{\fl}}\cong\rho_{2,\tilde{\fl}}$ by Proposition \ref{Main1}. Then the same Tate isogeny theorem argument as in the proof of Theorem \ref{mainthm1} follows.



\end{proof}

\section{Function Field Analogy of Strong Multiplicity One}

\subsection{Notation}

\begin{itemize}
\item $K:$ Global function field
\item $\Sigma_K:$ set of places of $K$, 
\item Let $S\subset \Sigma_K$ be a set of places, the Dirichlet density of S (if exists) is defined by
$$\delta(S):=\lim_{s\rightarrow 1^+}\frac{\sum_{\fp\in S}N(\fp)^{-s}}{\sum_{\fp\in \Sigma_K}N(\fp)^{-s}}=\lim_{s\rightarrow 1^+}\frac{\sum_{\fp\in S}N(\fp)^{-s}}{\log \frac{1}{s-1}}=\lim_{x\rightarrow \infty}\frac{\#\{\fp\in S\mid N(\fp)\leqslant x\}}{\#\{\fp\mid N(\fp)\leqslant x\}},$$
where $N(\fp):=|\cO_K/\fp|$ is the ideal norm.
\end{itemize}


\subsection{Algebraic Chebotarev density}
We start with the well known Chebotarev density theorem for finite extension of global fields

\begin{thm}[\cite{FJ08} , Section 6.4]\label{che}
Let $L$ be a finite Galois extension over $K$, and $\mathcal{C}$ be a conjugacy class of $G:=\Gal(L/K)$, then 
$$ \delta(\mathcal{C})=\frac{|\mathcal{C}|}{|G|}$$
\end{thm}

Now we generalize the Chebotarev density theorem based on the so-called local thinness property.

\begin{thm}\label{thma}

Let $E$ be a non-archimidean local field of positive characteristic $p$, and $\rho: G_K\rightarrow M(E)$ be a continuous representation that is unramified outside of a finite set of places in $\Sigma_K$, where $M$ is an affine algebraic group over $E$.  Let $H=\overline{\rho(G_K)}^{\rm Zar}\subset M$ be the Zariski closure of $\rho(G_K)$ in $M$, with identity component $H^\circ$, and finite quotient $\Phi:=H/H^\circ$.  Let $X\subset M$ be a conjugation-stable Zariski closed subset defined over $E$. Define
$$\Sigma_X=\left\{v\in\Sigma_K\mid \rho \textrm{ unramified at } v,\ \rho({\rm Frob}_v)\in X(E) \right\} {\rm and }\ \Psi=\left\{\phi\in\Phi\mid H^{\phi}\subset X \right\},$$ where $H^\phi$ denotes the corresponding connected component of $H$.

Assume the following componentwise thinness property:
\begin{equation}\tag{CT}
\textrm{For all $\phi\not\in\Psi$}, X(E)\cap\rho(G_K)\cap H^\phi(E) \textrm{ has Haar measure $0$ in }\rho(G_K)\cap H^\phi(E).
\end{equation}
Then $\Sigma_X$ has density $\delta(\Sigma_X)=\frac{|\Psi|}{|\Phi|}$.

\end{thm}

\begin{proof}

Denote $G=\rho(G_K)$ and $G^\circ=G\cap H^\circ(E)$, then we have the finite embedding $G/G^\circ\hookrightarrow \Phi$. Let $L/K$ be the finite Galois extension cut out by $G_K\rightarrow G/G^\circ$, i.e. $\Gal(L/K)\cong G/G^\circ \stackrel{\rho}\hookrightarrow\Phi$ . Let $v\in \Sigma_K$ be a place where $\rho$ is unramified at $v$. Then $\rho({\rm Frob}_v)$ embeds to a connected component of $H$ containing $\rho({\rm Frob}_v)$. Define $$\Sigma_{\rm comp}=\left\{v\in\Sigma_K\ {\rm unramified }\mid \rho({\rm Frob}_v)\in\cup_{\phi\in\Psi}H^\phi(E)\right\}.$$
Chebotarev density theorem then implies $$\delta(\Sigma_{\rm comp})=\frac{|\Psi|}{|\Phi|}.$$


Now we write $X\cap H=\left(\cup_{\phi\in\Psi}H^\phi \right)\cup Y$, where $Y\subset H$ is Zariski-closed satisfying the property that $\textrm{For any }\phi\not\in\Psi,\ Y\cap H^\phi=X\cap H^\phi$ is a proper Zariski-closed subset of $H^\phi$. Then we obtain a decomposition $$\Sigma_X=\Sigma_{\rm comp}\sqcup \Sigma_{\rm thin},$$ where $\Sigma_{\rm thin}:=\left\{v\in \Sigma_K\ {\rm unramified}\mid \rho({\rm Frob}_v)\in Y(E)\right\}$. By the property (CT), $\Sigma_{\rm thin}$ has Haar measure $0$ in $G\cap H^\phi(E)$ for any $\phi\not\in\Psi$. Lemma \ref{lema} below shows that $\delta(\Sigma_{\rm thin})=0$, and so $\delta(\Sigma_X)=\frac{|\Psi|}{|\Phi|}+0$.

\end{proof}

\begin{lem}\label{lema}
Up to conjugation, one can write $\rho$ as a continuous representation $\rho:G_K\rightarrow \GL_r(\cO_E)$. Let $A\subset G=\rho(G_K)$ be a measurable subset such that $\mu_G(A)=0$, where $\mu_G$ is the Haar measure on $G$ with $\mu_G(G)=1$. Define $\Sigma_A=\left\{v\in\Sigma_K\ {\rm unramified }\mid \rho({\rm Frob}_v)\in A\right\}$, then $\delta(\Sigma_A)=0$.
\end{lem}
\begin{proof}

For $n\in \mathbb{Z}_{\geqslant 1}$, consider $\pi_n:G\rightarrow G_n:=G/G\cap (1+\mathfrak{m}^n M_r(\cO_E))$ where $\mathfrak{m}$ is the maximal ideal of $\cO_E$ and $M_r(\cO_E)$ is the set of $r\times r$ matrices over $\cO_E$. Let $A_n=\pi_n(A)\subset G_n$, then $A\subset\pi_n^{-1}(A_n)$, which implies $\Sigma_A\subset\Sigma_{\pi_n^{-1}(A_n)}$.

For $x\in A_n$, $\pi_n^{-1}(x)$ are cosets of an open normal subgroup of $G$ of equal Haar measure, we can deduce that $$\mu_G(\pi_n^{-1}(A_n))=\frac{|A_n|}{|G_n|}.$$ Hence we have the property that $\frac{|A_n|}{|G_n|}\rightarrow 0$ as $n\rightarrow \infty$, otherwise the intersection $\cap_n\pi_n^{-1}(A_n)\supset A$ would have positive measure, a contradiction.

On the other hand, consider $\rho_n=\pi_n\circ\rho:G_K\rightarrow G_n$ and let $K_n/K$ be the finite Galois extension cut out by $\rho_n$, then
$$\Sigma_{\pi_n^{-1}(A_n)}=\left\{v\in\Sigma_K {\rm unramified}\mid \rho_n({\rm Frob}_v)\in A_n \right\}.$$
Theorem \ref{che} implies $$\delta(\Sigma_{\pi_n^{-1}(A_n)})=\frac{|A_n|}{|G_n|}.$$
From the fact that $\delta(\Sigma_A)\leqslant \delta(\Sigma_{\pi_n^{-1}(A_n)})=\frac{|A_n|}{|G_n|}$, we deduce $\delta(\Sigma_A)\leqslant\lim_{n\rightarrow \infty}\frac{|A_n|}{|G_n|}=0$.

\end{proof}


\subsection{Positive Characteristic SMO}
The main goal of this subsection is the following theorem:

\begin{thm}\label{thmb}
Let $K$ be a global function field, and  $\rho_1,\rho_2:G_K\rightarrow \GL_n(E)$ be two semisimple Galois representations over a local field $E$ of positive characteristic. Denote that
\begin{itemize}
\item $\rho:=(\rho_1,\rho_2):G_K\rightarrow \GL_n(E)\times\GL_n(E).$
\item$ H:=\overline{\rho(G_K)}^{\rm Zar}\subset \GL_r\times\GL_r$ is the Zariski closure of $\rho(G_K)$ in the affine algebraic group $\GL_r\times\GL_r$ over $E$.
\item $\Phi:=H/H^\circ$, where $H^\circ$ is the identity component of $H$.
\item $S:=\{v\in\Sigma_K\mid \rho\ {\rm unramified\ at\ }v\}$
\item$X:=\left\{(g_1,g_2)\in \GL_r\times\GL_r\mid \chi_{\rm char}(g_1)=\chi_{\rm char}(g_2) \right\}$, where $\chi_{\rm char}(g)$ is the characteristic polynomial of $g$. 
\item $\Sigma_X:=\{v\in S\mid \rho({\rm Frob}_v)\in X\}$
\end{itemize}

Suppose $H^\circ\subset X$, then we have the followings:
\begin{enumerate}
\item[(i)] There is a finite extension $L/K$ such that $\rho_1|_{G_L}\cong\rho_2|_{G_L}$
\item[(ii)]  If $\rho_1$ is absolutely irreducible and $\overline{\rho_1(G_K)}^{\rm Zar}\subset \GL_r$ is connected, then there is a finite character $\chi:G_K\rightarrow E^*$ such that $\rho_2\cong\rho_1\otimes \chi$.
\end{enumerate}
\end{thm}


Before diving into the proof of the theorem, let us state and prove some of its related applications.

\begin{cor}\label{cor1}
Under the notation of Theorem \ref{thmb}, suppose the componentwise thinness property (CT) in Theorem \ref{thma} holds for the product representation $\rho$. We have $H^\circ\subset X$ whenever the dirichlet density ${\delta}(\Sigma_X)>1-1/|\Phi|$.
\end{cor}
\begin{proof}
Apply Theorem \ref{thma}, we have $\delta(\Sigma_X)=\frac{|\Psi|}{\Phi}$ where $\Psi=\{\phi\in\Phi\mid H^\phi\subset X\}$. The assumption $\delta(\Sigma_X)>1-1/|\Phi|$ would force all connected components $H^\phi\subset X$; hence $H^\circ\subset X$.
\end{proof}

\begin{rem}\label{rem1}
One should compare the above corollary with the proof of Theorem 2,(i) in \cite{Ra98} for the case $E$ is a local field of characteristic $0$. Rajan's strategy starts with the assumption that the upper density of $\Sigma_X$ is positive and $H_1=\overline{\rho_1(G_K)}^{\rm Zar}$ is connected, then there is a $\phi\in\Phi$ such that the corresponding connected component $J^\phi:=H^\phi\cap J$ lies in $X$, here $J$ is a maximal compact subgroup of $H(\mathbb{C})$. Using the fact that maximal compact subgroup of $GL_r(\mathbb{C})$ is a group of unitary matrices, he can deduce that there is an element  $(I_r,j)\in J^\phi$ such that $j$ belongs to a group of unitary matrices $U_r(\mathbb{C})$. By trace counting and the fact that the only unitary matrix with trace $r$ is the identity matrix, $(I_r,j)$ is forced to be $(I_r,I_r)$, so $H^\phi=H^\circ\subset X$.   

In our situation where $E$ is of positive characteristic, we only have the fact that any maximal compact subgroup of $GL_r(E)$ is conjugate to $\GL_r(\cO_E)$, hence Rajan's strategy does not work in positive characteristic case. 


\end{rem}

\begin{cor}\label{cor2}
Assume Theorem \ref{thmb},(2) holds, i.e. $\rho_2\cong\rho_1\otimes \chi$ for some finite character $\chi:G_K\rightarrow E^*$, $\rho_1$ is absolutely irreducible and $H_1:=\overline{\rho_1(G_K)}^{\rm Zar}\subset \GL_r$ is connected. Suppose further that the local thinness property (LT) below holds for $\rho_1$
\begin{equation}\tag{LT}
\textrm{Any proper Zariski-closed subset $Z$ of $H_1$ has Haar measure $\mu_{\rho_1(G_K)}(\rho_1(G_K)\cap Z)=0$}
\end{equation}
, then the upper density $$\overline{\delta}(\Sigma_X):=\limsup_{s\rightarrow 1^+}\frac{\sum_{\fp\in S}N(\fp)^{-s}}{\log \frac{1}{s-1}}>1/2$$ implies $\rho_1\cong\rho_2$.

\end{cor}

\begin{proof}

If $\chi$ is trivial, then there is nothing to prove. Assume $\chi$ to be a nontrivial character, for $v\in S$, set $\zeta_v:=\chi({\rm Frob}_v)$. Then for $v\in \Sigma_X$, we have $$\chi_{\rm char}(\rho_1({\rm Frob}_v))=\chi_{\rm char}(\zeta_v\rho_1({\rm Frob}_v)).$$

For a  $1\not=\xi\in{\rm Im}(\chi)$, and define $Z_\xi:=\left\{g\in H_1\mid \chi_{\rm char}(g)=\chi_{\rm char}(\xi g)\right\}$. $Z_\xi$ is a conjugation-stable proper Zariski-closed subset of $H_1$ because $I_r\not\in Z_\xi$. Then the property (LT) implies $Z_\xi\cap \rho_1(G_K)$ has Haar measure 0 in $\rho_1(G_K)$. Hence Lemma \ref{lema} implies that $$\Sigma_\xi:=\{v\in S\mid \rho_1({\rm Frob}_v)\in Z_\xi(E)\}$$ has density $\delta(\Sigma_\xi)=0$.  Thus $\Sigma_{X,\xi}:=\{v\in \Sigma_X\mid \chi({\rm Frob}_v)=\xi\neq1\}$ has density equal to $0$, and we get $$\overline{\delta}(\Sigma_X)=\overline{\delta}\left(\Sigma_{X,1}\sqcup\sqcup_{1\neq\xi\in {\rm Im}\chi}\Sigma_{X,\xi}\right)=\overline{\delta}(\Sigma_{X,1})+\sum_{1\neq\xi\in{\rm Im}\chi}\overline{\delta}(\Sigma_{X,\xi})=\overline{\delta}(\Sigma_{X,1}),$$ 
here $\Sigma_{X,1}=\{v\in \Sigma_X\mid \chi({\rm Frob}_v)=1\}$. Moreover, by Theorem \ref{che} we have
$$\overline{\delta}(\Sigma_{X,1})\leqslant\overline{\delta}()\{v\in\Sigma_K\mid \chi({\rm Frob}_v=1)\}=\frac{1}{|{\rm Im}\chi|}\leqslant1/2,$$
which contradicts to the assumption $\overline{\delta}(\Sigma_X)>1/2$.


\end{proof}
\begin{rem}\label{rem2}
Apply Corollary \ref{cor1}, \ref{cor2} to $\fl$-adic Galois representations $\rho_1, \rho_2$ associated to rank-$r$ non-CM Drinfeld modules of generic characteristic. Since we have the adelic openness for each of $\rho_1$ and $\rho_2$ by the work of Pink-R\"utsche \cite{PR09}, we can say that
\begin{center}
Assuming property (CT) for product representation $\rho=(\rho_1,\rho_2)$, property (LT) for $\rho_1$, absolute irreducibility of $\rho_1$, then ${\delta}(\Sigma_X)> {\rm max}\left(1-\frac{1}{|\Phi|},\frac{1}{2}\right)$ implies $\rho_1\cong\rho_2$.
\end{center}

Note that one can actually derive property (LT) from the adelic openness of $\rho_1$, so what one should check is the property (CT) for product representation $\rho$. Besides, one can compare this result with Theorem 1 in \cite{Ra98} , for the characteristic $0$ phenomenon.  
\end{rem}



Finally, we begin the proof of our main goal.
\begin{proof}[proof of Theorem \ref{thmb}]

For $i\in\{1,2\}$, let $p_i:H\rightarrow GL_r$ be the projection onto the $i$-th coordinate. As $H^\circ\subset X$, we have 
$$\chi_{\rm char}(p_1(h))=\chi_{\rm char}(p_2(h))\ \textrm{for all }h\in H^\circ.$$
Now because $\rho_1$ and $\rho_2$ are semisimple  representations of $G_K$, we have that $p_1|_{H^\circ}$ and $p_2|_{H^\circ}$ are semisimple representations of $H^\circ$. Theorem \ref{BN} then implies 
\begin{equation}\tag{1}
p_1|_{H^\circ}\cong p_2|_{H^\circ}
\end{equation}
Let $\pi:G_K\xrightarrow{\rho}H(E)\twoheadrightarrow \Phi$ and set $L/K$ be the finite extension cut out by ${\rm ker}(\pi)$, then $\rho(G_L)\subset H^\circ(E)$ and (1) restrict further on $\rho(G_L)$ gives 
\begin{equation}\tag{2}
\rho_1|_{G_L}\cong\rho_2|_{G_L},
\end{equation}
this proves (i).

One may replace $L$ by its Galois closure in $G_K$ while (2) still holds. Thus, without loss of generality, we may assume $L/K$ is Galois.

Now assume $\rho_1$ is absolutely irreducible and $H_1:=\overline{\rho_1(G_K)}^{\rm Zar}\subset \GL_r$ is connected. Let $V_i$ be the $E$-space of $\rho_i$ for $i=1,2$, choose a $G_L$-equivariant isomorphism $f:V_1\rightarrow V_2$. 

\begin{claim*}
${\rm dim}_E({\rm Hom}_{G_L}(V_1,V_2))=1$
\end{claim*}

One the claim holds, for $\sigma\in G_K$ we define $f_\sigma:=\rho_2(\sigma)f\rho_1(\sigma)^{-1}$. As $G_L\lhd G_K$, for $h\in G_L$ one can deduce that
$$h\cdot f_\sigma=\rho_2(\sigma)\rho_2(\sigma^{-1}h\sigma)f\rho_1(\sigma)^{-1}=\rho_2(\sigma)f\rho_1(\sigma^{-1}h\sigma)\rho_1(\sigma)^{-1}=f_\sigma\cdot h.$$
Thus $f_\sigma\in {\rm Hom}(V_1,V_2)$. By the claim, ${\rm Hom}(V_1,V_2)$ is the $E$-space spanned by $f$, hence there is an element $\chi(\sigma)\in E^*$ such that $f_\sigma=\chi(\sigma)f$. This defines a map
$$\chi:G_K\rightarrow E^*.$$
For $\sigma_1,\sigma_2\in G_K$, we check that 
$$\chi(\sigma_1\sigma_2)f=f_{\sigma_1\sigma_2}=\rho_2(\sigma_1)f_2\rho_1(\sigma_1)^{-1}=\chi(\sigma_2)f_{\sigma_1}=\chi(\sigma_1)\chi(\sigma_2)f.$$
Hence $\chi$ is a group homomorphism. Together with the fact that $\chi$ is continuous since $\sigma\mapsto f_\sigma$ is continuous, we know that $\chi$ is a character of $G_K$.  Moreover, from the property that $f_h=f$ for any $h\in G_L$, we know $\chi$ factors through the finite group ${\rm Gal}(L/K)$, so $\chi$ is a finite character.
Now from $f_\sigma=\chi(\sigma)f$ we get 
$$\rho_2(\sigma)f=f\chi(\sigma)\rho_1(\sigma)=f((\rho_1\otimes\chi)(\sigma)).$$
This means $f$ is a $G_K$-module isomorphism between $(V_1,\rho_1\otimes\chi)$ and $(V_2,\rho_2)$, hence $\rho_2\cong\rho_1\otimes\chi$.

\begin{proof}[proof of claim]

Choose a $G_L$-equivariant isomorphism $f:V_1\rightarrow V_2$. Consider the map
$$\eta:{\rm End}_{G_L}(V_1)\rightarrow {\rm Hom}_{G_L}(V_1,V_2),\ \ u\mapsto f\circ u.$$
One can easily see that $\eta$ is an $E$-isomorphism. Thus we turn to prove ${\rm dim}_E({\rm End}_{G_L}(V_1))=1$.

Since $H_1$ is connected, the finite index subgroup $\rho_1(G_L)\subset \rho_1(G_K)$ is still Zariski dense in $H_1$. Hence one can deduce that ${\rm End}_{G_L}(V_1)={\rm End}_{H_1}(V_1)={\rm End}_{G_K}(V_1)$, then Schur's lemma implies
$${\rm End}_{G_L}(V_1)={\rm End}_{G_K}(V_1)=E.$$
The proof is complete
\end{proof}
\end{proof}



\bibliographystyle{alpha}
\bibliography{Frobenius_trace_DM_rk2.bib}

\end{document}
